\subsection{Course and Curriculum}
\label{sec:impl_course}

Courses are the unit a learner enrols in and the container every other feature hangs from. The implementation lives in \path{abridgeai/features/courses/}, split by audience rather than by entity: \path{authoring.py} for the teacher's editing surface, \path{catalog.py} for the learner's browse and detail views, \path{assignment.py} for teacher--course bindings, and \path{administration.py} for the operator actions that cross both.

\paragraph{The curriculum tree}
A course owns ordered modules, and each module owns ordered items --- lessons, quizzes, and interview configurations addressed uniformly so the learner's next step is a single ordered walk rather than three parallel lists. Ordering is a dense \texttt{position} integer per parent under a unique constraint, so two items can never claim the same slot; reordering rewrites the affected range in one statement rather than issuing a move per row.

\paragraph{Publication gating}
A course, like every other authored artefact in the system, moves through \texttt{draft}, \texttt{published}, and \texttt{archived}. Draft content is invisible to learners at the query level rather than filtered in the presentation layer, so an unpublished lesson cannot leak through an endpoint that forgot to check.

Neither transition is reversible, and the asymmetry is deliberate. Publication is a one-way door: a published course can never return to draft, because its learning outcomes double as the graded assessment scale and re-opening it for editing would move the goalposts under students already enrolled against them. Archival is the separate terminal state, reached from either of the others, and it withdraws a course from new enrolment without disturbing the enrolments already running through it. A course removed in error is recovered not by a state change but by restoring the soft-deleted row, which is an operator action with its own audit trail rather than an ordinary edit.

\paragraph{Syllabus import}
\path{syllabus_import.py} accepts a structured syllabus and materialises an entire module tree in one operation. It follows the same per-row discipline as the programme roster import: a malformed entry is reported against its position and the remainder still imports, because a syllabus with one bad row is the normal case rather than the exception.

\paragraph{Enrolment and entitlement}
Learners reach a course three ways --- operator enrolment, a career path or programme that grants it, and self-enrolment by invitation code --- and all three converge on one enrolment row whose status is \texttt{active}, \texttt{completed}, \texttt{dropped}, or \texttt{waitlisted}. Keeping the grant paths separate from the enrolment record is what allows a programme path change to release an entitlement without disturbing an enrolment the learner obtained independently.
